\documentclass[12pt]{article} % Тип документа (стаття) та розмір шрифту (12pt)
\usepackage[T2A]{fontenc}      % Кодування шрифтів для відображення кирилиці
\usepackage[utf8]{inputenc}    % Кодування вхідного файлу (стандарт для укр.мови)
\usepackage[ukrainian]{babel}  % Підключення словників та правил переносу укр.мови
\usepackage{amsmath, amssymb}  % Пакети для математичних формул та символів
\usepackage{array}             % Пакет для покращеної роботи з таблицями

\title{Практична робота: Математичні формули} 
\author{Бондар Наталія}                          
\date{\today}                                

\begin{document} % Початок тексту документа

\maketitle % Команда, яка фактично малює заголовок, автора та дату на сторінці
\section*{Квадратне рівняння}
\begin{equation}
ax^2 + bx +c = 0,\qual a\neq 0
\end{equation}
\begin{table}[h]
\centering
\renewcommand{\arraystretch}{3}
\begin{tabular}{|c|c|c|}
\hline
\multicolumn{3}{|c|}{$ax^2 + bx + c = 0, a \neq 0, b \neq 0, c \neq 0$}\\ \hline
\multicolumn{3}{|c|}{$D = b^2 - 4ac$}\\ \hline
$\begin{aligned}
    x_1 &= \fract{-b + \sqrt{D}}{2a} \\
    x_2 &= \fract{-b - \sqrt{D}}{2a}
\end{aligned}$ &
$x = -\dfrac{b}{2a}$ &


\end{tabular}
\end{table}
\end{document} % Кінець документа
